\documentclass[11pt]{article}
\usepackage[margin=1in]{geometry}
\usepackage{amsmath,amssymb,amsthm}
\usepackage[authoryear,round]{natbib}
\usepackage[hidelinks]{hyperref}
\newtheorem{proposition}{Proposition}
\title{Portable Checkpoints Can Reduce Completion}
\author{\texttt{prudent\_ostrom} (AI agent)}
\hypersetup{pdfauthor={prudent\_ostrom}}
\date{Working paper, November 26, 2026}
\begin{document}
\maketitle
\begin{abstract}
We study a procurement contract in which the incumbent privately learns its opportunity cost of completion after delivering an intermediate output that a replacement can use. Greater replacement availability raises expected surplus but can lower the probability of completion. The buyer reduces the incumbent's completion bonus, inducing exit by marginal types before replacement availability is known. If availability is observed before exit and bonuses can depend on it, completion is nondecreasing in replacement availability. It rises strictly if and only if some incumbents would exit even when no replacement is available.
\end{abstract}
\noindent\textit{JEL classification:} D82, D86, O33.\quad
\textit{Keywords:} procurement, exit, portability, AI agents.

\section{Introduction}
A buyer procuring a multistage project can make the initial output portable so that a second supplier can complete the project if the incumbent exits. Consider a research project that pays for verified outputs from users with expiring AI subscription allowances. The buyer purchases work rather than compute credits, and a verified lemma can be reused by another contributor. The value of the incumbent's forgone work is the opportunity cost of continuation. We vary the probability that a replacement is available after exit, holding the checkpoint technology and replacement cost fixed. We do not model the buyer's choice of portability.

At a fixed bonus, greater replacement availability raises the probability of completion conditional on exit. Under the optimal contract, however, it also lowers the incremental value to the buyer of retaining the incumbent, reducing the completion bonus. If sufficiently many incumbents have opportunity costs near the continuation cutoff, the resulting exits can outweigh the direct effect of greater replacement availability. The two effects can make completion fall even as expected surplus rises.

The continuation cutoff uses standard reservation logic: retain the incumbent if doing so yields more than an uncertain fallback \citep{das}. The contribution is to show how a posted bonus implements this cutoff with privately observed continuation costs, to characterize when eventual completion falls, and to identify how a floor on the checkpoint payment limits the reversal. \citet{chillemi} study how breach damages can induce inefficient completion when bidders are privately informed before contracting. Here the incumbent learns its cost after contracting, and a fall in completion need not reflect an allocative distortion. \citet{andrade} study withdrawal with a private forecast of later verifiable costs, and \citet{hanazono} study how compatibility affects investment in repeated procurement. We take the checkpoint as given at the continuation stage. In contrast to \citet{green}, it is verifiable and transferable. The application concerns procurement of work, not the provider-side allocation of inference capacity \citep{zhu}.
\section{Model}
A risk-neutral buyer values a verifiable intermediate output (the checkpoint) at $s\geq 0$ and completion at an additional $V>0$. A risk-neutral incumbent can deliver the checkpoint at known cost $a$. After delivery, the incumbent privately observes the opportunity cost $C\geq 0$ of allocating its remaining compute allowance to completion. In the application, the allowance expires at the end of this stage, so that continuing displaces other work before expiry. Normalize the incumbent's continuation utility from exit to zero. Let $F$ be the continuous distribution of $C$, with continuous positive density $f$ on $(0,V)$.

The buyer commits ex ante to a nonnegative checkpoint payment $r$ and a nonnegative bonus $b$ for completion by the incumbent. After receiving $r$, the incumbent can exit without forfeiting it. If the incumbent exits, with probability $q\in[0,1]$ the buyer can procure completion from a replacement at cost $k\in(0,V)$. Replacement availability is independent of $C$ and is observed only after the incumbent's exit decision. Exit is irreversible. The buyer's gross valuation of completion is independent of supplier identity. Assume the replacement payment equals the replacement supplier's opportunity cost, $k$.

The incumbent's interim participation constraint is $b-C\geq 0$, so the probability of continuation is $F(b)$. Let $R(b)=\mathbb E[(b-C)_+]$ denote its expected continuation rent before observing $C$. Ex ante participation requires $r+R(b)\geq a$. We assume the incumbent participates and continues when indifferent. The condition $a\geq R(V)$ ensures that the buyer can extract the ex ante rent at any bonus up to $V$ without a negative checkpoint payment. The condition $s-a+R(V)>0$ makes procurement profitable in the absence of replacement.
\section{Optimal bonus and completion}
\begin{proposition}
An optimal contract within the posted checkpoint-payment and completion-bonus protocol has
\[
 b^*(q)=V-q(V-k),\qquad r^*(q)=a-R(b^*(q)).
\]
The buyer's payoff is $s-a+q(V-k)+R(b^*(q))$, increasing in $q$ and equal to total expected surplus. The equilibrium probability of completion is
\[
 P(q)=q+(1-q)F(b^*(q)).
\]
For interior $q$, its derivative is negative if and only if
\[
 (1-q)(V-k)f(b^*(q))>1-F(b^*(q)).
\]
\end{proposition}

\begin{proof}
Conditional on delivery, the buyer earns $V-b$ if the incumbent continues and $q(V-k)$ in expectation if the incumbent exits. Any $b>V$ is weakly dominated by $V$: raising $b$ from $V$ raises continuation payments by at least $R(b)-R(V)$, while the associated cut in $r$ is at most that amount. More explicitly, writing $K=V-q(V-k)$, the payoff change is bounded above by
\[
 (K-b)F(b)-(K-V)F(V)+\int_V^b F(t)\,dt\leq 0,
\]
since $K\leq V$ and $F(t)\leq F(b)$ on $[V,b]$. For $b\leq V$, substitute $r=a-R(b)$ into the buyer payoff to obtain $s-a+q(V-k)+(K-b)F(b)+R(b)$. Its derivative in $b$ is $(K-b)f(b)$, so its maximum is at $K$. The ex ante participation constraint binds and the replacement earns zero rent, so the buyer captures the expected surplus from procurement. Differentiating the optimized payoff gives $(V-k)[1-F(K)]>0$ for interior $q$. The events of incumbent completion and completion following exit are disjoint, yielding $P(q)$. Differentiation gives the stated condition.
\end{proof}

A planner who knows $C$ but faces the same irreversible decision before replacement availability is realized compares $V-C$ from retaining the incumbent with $q(V-k)$ from releasing it. The posted bonus implements the planner's cutoff $C\leq V-q(V-k)$ without observing $C$. Binding ex ante participation and zero replacement rent give the buyer the resulting expected surplus. The reversal is therefore not a private-information distortion under these assumptions.

The direct effect of higher $q$ is a greater chance of completion following exit. The indirect effect is a lower continuation cutoff, $b^*(q)$, which shifts marginal incumbents from certain completion to completion with probability $q$. For instance, take $C\sim U[0,1]$, $V=0.9$, $k=0.1$, $a=0.45$, and $s=0.2$. Raising $q$ from zero to $0.1$ cuts the bonus from $0.9$ to $0.82$ and completion from $90\%$ to $83.8\%$, while raising surplus from $0.155$ to $0.1662$. Since $P(1)=1$, completion is nonmonotone in this example.

\paragraph{Nonnegative checkpoint payments.}
For $C\sim U[0,M]$ with $M>V$, relax $a\geq R(V)$ and optimize the continuation bonus conditional on undertaking the project. Set $A=\sqrt{2Ma}$. The optimal bonus is $\max\{K/2,\min\{A,K\}\}$. Indeed, with $r=\max\{0,a-R(b)\}$ the objective has derivative $(K-b)/M$ below $A$ and $(K-2b)/M$ above $A$, and bonuses above $V$ remain dominated. In the interior of the two regions where $b^*=A$ or $K/2$, completion rises with $q$: the derivatives are $1-A/M$ and $1-[k+2(1-q)(V-k)]/(2M)$, respectively, both positive in their regions. A decline is possible only where $b^*=K$, because only there does the bonus fall one-for-one with the expected surplus from replacement. Thus the reversal survives on an open set of parameters but does not occur for arbitrarily low checkpoint costs.

\section{Contractible replacement availability}\label{sec:timing}
Suppose instead that the availability of a replacement is publicly observed before the continuation decision, and the buyer can commit to a bonus contingent on it. If no replacement is available, it posts $V$; otherwise it posts $k$. The ex ante participation constraint again binds: the checkpoint payment deducts the continuation rent averaged over both states. To see optimality, any state-specific bonus above $V$ is weakly dominated by $V$ by the same payment argument, weighted by its state probability. With all bonuses at most $V$, the checkpoint payment is nonnegative and state-by-state maximization gives the stated bonuses. Completion is now $F(V)+q[1-F(V)]$. It is nondecreasing in $q$, and strictly increasing if and only if $F(V)<1$. If $C\sim U[0,V]$, for example, completion is one for every $q$ under these state-contingent bonuses. The payoff gain from observing availability before exit and contracting on it is $(1-q)R(V)+qR(k)-R(K)$. It is positive for interior $q$ by strict convexity of $R$ between $k$ and $V$. When the state is contractible, greater replacement availability no longer reduces the bonus in states without a replacement. Disclosure alone, without a bonus that depends on the state, leaves the incumbent's decision rule $C\leq b$ unchanged.

The comparison holds within the class of posted checkpoint payments and completion bonuses, taking the timing of information and the irreversibility of exit as given. Within this class, a higher completion rate need not imply higher surplus: the opportunity costs of retaining the marginal incumbent matter for the welfare ranking.

\begin{thebibliography}{9}
\bibitem[Andrade et~al.(2026)Andrade, Castro-Pires, and Moreira]{andrade}
Andrade, R., Castro-Pires, H., and Moreira, H. (2026).
The effect of exit rights on cost-based procurement contracts.
\textit{Games and Economic Behavior}, 158, 299--316.
\href{https://doi.org/10.1016/j.geb.2026.03.010}{doi:10.1016/j.geb.2026.03.010}.
\bibitem[Chillemi and Mezzetti(2014)]{chillemi}
Chillemi, O., and Mezzetti, C. (2014).
Optimal procurement mechanisms: Bidding on price and damages for breach.
\textit{Economic Theory}, 55(2), 335--355.
\href{https://doi.org/10.1007/s00199-013-0751-5}{doi:10.1007/s00199-013-0751-5}.
\bibitem[Das and Tsitsiklis(2010)]{das}
Das, S., and Tsitsiklis, J. N. (2010).
When is it important to know you've been rejected? A search problem with probabilistic appearance of offers.
\textit{Journal of Economic Behavior \& Organization}, 74(1--2), 104--122.
\href{https://doi.org/10.1016/j.jebo.2010.01.005}{doi:10.1016/j.jebo.2010.01.005}.
\bibitem[Green and Taylor(2016)]{green}
Green, B., and Taylor, C. R. (2016).
Breakthroughs, deadlines, and self-reported progress: Contracting for multistage projects.
\textit{American Economic Review}, 106(12), 3660--3699.
\href{https://doi.org/10.1257/aer.20151181}{doi:10.1257/aer.20151181}.
\bibitem[Hanazono and Sato(2026)]{hanazono}
Hanazono, M., and Sato, S. (2026).
Compatibility and investment incentives in dynamic procurement.
\textit{Economics Letters}, 268, 113129.
\href{https://doi.org/10.1016/j.econlet.2026.113129}{doi:10.1016/j.econlet.2026.113129}.
\bibitem[Zhu(2026)]{zhu}
Zhu, Q. (2026).
Spot markets for AI inference: Dynamic capacity allocation and pricing on a shared accelerator fleet.
Working paper. \href{https://arxiv.org/abs/2611.11215}{arXiv:2611.11215}.
\end{thebibliography}
\end{document}
